% Part: first-order-logic
% Chapter: beyond
% Section: second-order-logic

\documentclass[../../../include/open-logic-section]{subfiles}

\begin{document}

\olfileid{fol}{byd}{sol}

\olsection{二阶逻辑}

二阶逻辑的语言不仅允许对个体论域进行量化，还允许对该论域上的关系进行量化。给定一阶语言 $\Lang{L}$，对于每个 $k$，增加取值遍历 $k$ 元关系的变元 $R$，并允许对这些变元进行量化。如果 $R$ 是 $k$ 元关系变元，且 $t_1$, \dots,~$t_k$ 是普通（一阶）项，则 $\Atom{R}{t_1,\dots,t_k}$ 是原子公式。除此之外，公式集合的定义与一阶逻辑的情形相同，只需为二阶量化添加相应子句。注意，我们只有关于一阶项的等词谓词：若 $R$ 和 $S$ 是相同元数 $k$ 的关系变元，我们可以将 $\eq[R][S]$ 定义为
\[
\lforall[x_1 \dots][\lforall[x_k][(\Atom{R}{x_1, \dots, x_k} \liff
  \Atom{S}{x_1, \dots, x_k})]].
\]
的缩写。

二阶逻辑的规则只需将量词规则推广到新的二阶变元。然而，在此必须稍加小心，以解释这些变元如何与 $\Lang{L}$ 的谓词符号乃至更一般地与 $\Lang{L}$ 的公式相互作用。最起码，关系变元被视作项，因此我们有如下形式的推理：
\[
!A(R) \Proves \lexists[R][!A(R)]
\]
但是，如果 $\Lang{L}$ 是带有常关系符号 $<$ 的算术语言，我们也会期望下列推理有效：
\[
x < y \Proves \lexists[R][\Atom{R}{x,y}]
\]
或者，对于给定的公式 $!A$，
\[
!A(x_1, \dots, x_k) \Proves \lexists[R][\Atom{R}{x_1,\dots,x_k}]
\]
更一般地，我们可能想要允许如下形式的推理：
\[
\Subst{!A}{\lambd[\vec x][!B(\vec x)]}{R} \Proves
\lexists[R][!A]
\]
其中 $\Subst{!A}{\lambd[\vec x][!B(\vec x)]}{R}$ 表示将 $!A$ 中所有形如
$\Obj{R}{t_1,\dots,t_k}$ 的原子公式替换为 $!B(t_1, \dots, t_k)$ 后得到的结果。最后这条规则等价于拥有一个{\em 概括图式}（comprehension schema），即对二阶语言中的每个公式 $!A$ 都有一条形如
\[
\lexists[R][\lforall[x_1, \dots, x_k][(!A(x_1, \dots, x_k) \liff
\Atom{R}{x_1, \dots, x_k})]],
\]
的公理，且 $R$ 不在 $!A$ 中自由出现。（练习：证明如果允许 $R$ 在 $!A$ 中出现，则该图式是不一致的！）

当逻辑学家提及“二阶逻辑的公理”时，通常指的是一阶逻辑经二阶量词规则连同概括图式的最小扩展。但研究这些公理与规则的较弱子系统通常也很有趣。例如，注意概括图式在其最一般的形式下是\emph{非直谓的}（impredicative）：它允许断言存在一个由带二阶量词的公式“定义”的关系 $\Atom{R}{x_1, \dots, x_k}$；而这些量词遍历所有此类关系的集合——该集合包含 $R$ 自身！在20世纪初，对罗素悖论的一种常见反应是将责任归咎于此类定义，并在发展数学基础时避免使用它们。若禁止在公式 $!A$ 中使用二阶量词，就得到概括图式的一种{\em 直谓}（predicative）形式，其强度稍弱。

从语义角度看，可以将二阶结构视为由语言的一阶结构，耦合上二阶量词所遍历的论域上的关系集合构成（更确切地说，对每个 $k$ 都有一个 $k$ 元关系的集合）。当然，若推导系统中包含概括图式，则需增加一个要求，即“二阶部分”中必须有足够的关系来满足概括公理——否则推导系统就不是可靠的！确保有足够关系的一个简单方法是让二阶部分由一阶部分上的\emph{所有}关系构成。这种结构被称为\emph{完全的}（full），并且在某种意义上，它确实是该语言的“预期结构”。如果我们只关注完全结构，就得到了所谓的\emph{完全}二阶语义。此时，指定一个结构就归结为指定其一阶部分，因为二阶部分的内容可由此隐含地确定。

总之，在谈论二阶逻辑时存在一些歧义。就推导系统而言，可能指的是：

\begin{enumerate}
\item 一个“最小的”二阶推导系统，加上一些概括公理。
\item “标准的”二阶推导系统，具有完全的概括图式。
\end{enumerate}

就语义而言，可能感兴趣的是：

\begin{enumerate}
\item “弱”语义，其中结构由一阶部分加上足以满足概括公理的二阶部分构成。
\item “标准的”二阶语义，其中只考虑完全结构。
\end{enumerate}

当逻辑学家未特别指明所考虑的推导系统或语义时，通常指的是上述两个列表中的第二项。使用这种语义的优势在于，正如我们将要看到的，它能为许多自然数学结构提供范畴描述；同时，推导系统相当强大，且对该语义是可靠的。其缺点在于，推导系统对该语义是\emph{不完备的}；事实上，\emph{没有}任何能行给出的推导系统能对完全二阶语义是完备的。另一方面，我们将看到，推导系统对弱化语义\emph{是}完备的；这意味着，若一个语句不可证，则存在\emph{某个}结构（不一定是完全结构），在其中该语句为假。

二阶逻辑的语言相当丰富。我们可以将一元关系等同于论域的子集，因此特别地，我们可以对这些集合进行量化；例如，可以用单条公理表达自然数的归纳法：
\[
\lforall[R][((\Atom{R}{\Obj{0}} \land \lforall[x][(\Atom{R}{x} \lif
    \Atom{R}{x'})]) \lif \lforall[x][\Atom{R}{x}])].
\]
若取算术语言包含符号 $\Obj 0, \Obj \prime, +, \times$ 和 $<$，则可添加以下公理来描述它们的行为：

\begin{enumerate}
\item $\lforall[x][\lnot x' = \Obj 0]$
\item $\lforall[x][\lforall[y][(s(x) = s(y) \lif x = y)]]$
\item $\lforall[x][(x + \Obj 0) = x]$
\item $\lforall[x][\lforall[y][(x + y') = (x + y)']]$
\item $\lforall[x][(x \times \Obj 0) = \Obj 0]$
\item $\lforall[x][\lforall[y][(x \times y') = ((x \times y) + x)]]$
\item $\lforall[x][\lforall[y][(x < y \liff \lexists[z][y = (x + z')])]]$
\end{enumerate}

不难证明，只要使用完全二阶语义，这些公理连同上面的归纳公理，就为结构 $\Struct{N}$（算术的标准模型）提供了范畴描述。给定任意满足这些公理的结构 $\Struct{M}$，利用 $\Nat$ 上的普通递归来定义一个从 $\Nat$ 到 $\Struct{M}$ 的论域的函数 $f$，使得 $f(0) = \Assign{\Obj 0}{M}$ 且 $f(x+1) = \Assign{\prime}{M}(f(x))$。利用 $\Nat$ 上的普通归纳法以及公理 (1) 和 (2) 在 $\Struct M$ 中成立的事实，可知 $f$ 是单射。为证明 $f$ 是满射，令 $P$ 为 $\Struct{M}$ 论域中属于 $f$ 值域的元素组成的集合。由于 $\Struct M$ 是完全结构，$P$ 就在其二阶论域中。根据 $f$ 的构造，可知 $\Assign{\Obj 0}{M}$ 属于 $P$，且 $P$ 在 $\Assign{\prime}{M}$ 下封闭。归纳公理在 $\Struct M$ 中成立这一事实（特别是对 $P$ 成立）保证了 $P$ 等于 $\Struct M$ 的整个一阶论域。这表明 $f$ 是双射。反复利用 $\Nat$ 上的普通归纳法，要证明 $f$ 是同态也不困难。

用集合论的术语来说，函数只是一种特殊的关系；例如，一元函数 $f$ 可等同于满足 $\lforall[x][\lexists![y][R(x,y)]]$ 的二元关系 $R$。因此，也可以对函数进行量化。利用完全语义，可以将无穷结构的类定义为满足如下条件的结构 $\Struct M$ 的类：存在从 $\Struct M$ 的论域到其真子集的单射函数：
\[
\lexists[f][(\lforall[x][\lforall[y][(\eq[f(x)][f(y)] \lif \eq[x][y])]]
  \land \lexists[y][\lforall[x][\eq/[f(x)][y]]])].
\]
该语句的否定则定义了有穷结构的类。

此外，可以通过在线性序的定义上添加如下条件来定义良序的类：
\[
\lforall[P][(\lexists[x][\Atom{P}{x}] \lif \lexists[x][(\Atom{P}{x}
    \land \lforall[y][(y < x \lif \lnot \Atom{P}{y})])])].
\]
这断言了每个非空“集合”都有最小元，只需将“集合”等同于“一元关系”。再举一例，我们可以表达图的连通性概念，即不存在将顶点非平凡地分割成不连通部分的划分：
\[
\lnot \lexists[A][(\lexists[x][A(x)] \land \lexists[y][\lnot A(y)]
  \land \lforall[w][\lforall[z][((\Atom{A}{w} \land \lnot \Atom{A}{z})
      \lif \lnot \Atom{R}{w,z})]])].
\]
作为练习，读者可以尝试定义论域大小为偶数的有穷结构的类。更引人注目的是，我们可以将实数范畴地描述为包含有理数的完备有序域。

简言之，二阶逻辑比一阶逻辑的表达能力强得多。这是好消息；现在来说坏消息。前文已提及，对于完全二阶语义，不存在完备的能行推导系统。无论好坏，一阶逻辑的许多性质在二阶逻辑中都不复存在，包括紧致性和 L\"owenheim--Skolem（勒文海姆-斯科伦）定理。

另一方面，如果愿意用较弱的语义代替完整的二阶语义，那么极小二阶推导系统相对于这种语义就是完备的。换句话说，如果我们将 $\Proves$ 理解为“在极小系统中可推导”，将 $\Entails$ 理解为“在较弱语义下语义地后承”，我们就能证明：若 $\Gamma \Entails !A$，则 $\Gamma \Proves !A$。如果要在推导系统中加入特定的概括公理，就必须将语义限制在满足这些公理的二阶结构上：例如，若 $\Delta$ 由一组（可能是全部）概括公理组成，则若 $\Gamma \cup \Delta \Entails !A$，就有 $\Gamma \cup \Delta \Proves !A$。特别地，若利用我们所考虑的概括公理不能证明 $!A$，那么就存在 $\lnot !A$ 的一个模型，其中这些概括公理仍然成立。

理解较弱语义下完备性定理成立的最简单方式，是将二阶逻辑视为一种多类逻辑，具体如下。一个类被解释为通常的“一阶”论域，而对于每个 $k$，我们另有一个“$k$ 元关系的论域”。我们取语言中包含内置的关系符号“$\Atom{\Obj{true}_k}{R,x_1,\dots,x_k}$”，它意在断言关系 $R$ 对 $x_1$, \dots,~$x_k$ 成立，其中 $R$ 是“$k$ 元关系”类的一个变元，而 $x_1$, \dots,~$x_k$ 是一阶类的对象。

在这种等同下，弱二阶语义本质上就是多类逻辑的通常语义；而我们已知，多类逻辑可以嵌入到一阶逻辑中。因此，在双向翻译的意义上，这种较弱的二阶逻辑观念实际上就是一阶逻辑的一种伪装形式，其论域同时包含“对象”和“关系”，并受相应公理的约束。

\end{document}